\section{Назначение системы}

Проект реализует полный IoT-контур для мониторинга параметров окружающей среды в реальном времени. Система предназначена для непрерывного сбора телеметрии, визуализации текущих значений, накопления исторических данных и поддержки удалённого контроля состояния объекта наблюдения.

С инженерной точки зрения решение закрывает четыре ключевые задачи:

\begin{enumerate}
    \item \textbf{Сбор и первичная обработка телеметрии.} Микроконтроллер (ESP8266 / NodeMCU v2) считывает показания датчиков температуры, влажности и концентрации газа, выполняет базовую очистку данных и локально отображает значения на OLED-дисплее.
    \item \textbf{Frontend (React + TypeScript)} читает данные из Firebase:
    \begin{itemize}
        \item \texttt{current} --- «живые» текущие показания;
        \item \texttt{history} --- почасовая история.
    \end{itemize}
    \item \textbf{Оптимизация трафика и нагрузки.} Frontend отправляет heartbeat (\texttt{/status/last\_viewed}), по которому прошивка понимает, активен ли пользователь, и решает, публиковать ли ветку \texttt{current}.
    \item \textbf{Формирование аналитической базы.} Прошивка сохраняет агрегированную историю в Firebase по часовым срезам, что позволяет строить графики, анализировать тренды и оценивать динамику параметров без постоянной потоковой передачи.
\end{enumerate}

\subsection{Функциональные границы решения}

В текущей реализации система ориентирована на телеметрию и визуализацию, без подсистемы управления исполнительными устройствами. Это сознательный выбор архитектуры: контур «измерение → хранение → отображение» изолирован от контура «управление → обратная связь», что снижает риск ошибок и упрощает сопровождение.

\begin{itemize}
    \item \textbf{Входит в область решения:} сбор метрик, публикация актуальных данных, хранение истории, интерактивный просмотр в веб-интерфейсе.
    \item \textbf{Не входит в область решения:} аварийные команды на оборудование, пользовательские роли доступа, локальное долговременное хранилище на устройстве.
\end{itemize}

\subsection{Критерии качества отчётной системы}

Система оценивается по следующим базовым критериям:

\begin{itemize}
    \item \textbf{Надёжность сбора данных:} цикл измерений выполняется периодически и не блокируется сетевыми операциями.
    \item \textbf{Согласованность времени:} все временные метки приводятся к UTC epoch, что исключает ошибки при интерпретации данных на разных клиентах.
    \item \textbf{Читаемость интерфейса:} пользователь получает как текущую картину, так и исторический контекст в удобном формате.
    \item \textbf{Масштабируемость по данным:} архитектура разделяет «live»-ветку и историческую ветку, снижая нагрузку при росте объёма измерений.
\end{itemize}

\section{Состав репозитория}

\subsection{Логическая декомпозиция проекта}

\begin{itemize}
    \item \texttt{course\_work/} --- прошивка микроконтроллера (PlatformIO, Arduino framework). Здесь сосредоточены аппаратная инициализация, циклический сбор данных, синхронизация времени и обмен с Firebase.
    \item \texttt{iot-dashboard/} --- веб-панель мониторинга (Vite + React + TypeScript + MUI + Recharts + Firebase SDK). Модуль отвечает за отображение текущих и исторических данных, пользовательские фильтры и heartbeat-активность.
    \item \texttt{generate/} --- вспомогательные скрипты генерации данных (не участвуют в основном runtime-контуре прошивка↔frontend). Используются для подготовки тестовых/демо-наборов и отладки визуализаций.
\end{itemize}

\subsection{Принцип разделения ответственности}

Разделение по директориям выполнено в соответствии с принципом единой ответственности: микроконтроллерный код не содержит UI-логики, а frontend не содержит аппаратно-зависимых алгоритмов. Это позволяет независимо развивать обе части системы и упрощает тестирование.

\section{Архитектура на уровне подсистем}

На диаграмме показан полный путь телеметрии: от физического измерения до отображения на клиенте. Центральная роль у Firebase RTDB как у точки синхронизации между асинхронной прошивкой и реактивным веб-интерфейсом.

\begin{itemize}
    \item Подсистема ESP8266 формирует данные и инициирует запись.
    \item Firebase хранит состояние и обеспечивает подписки/чтение.
    \item Frontend интерпретирует данные, отображает их и формирует heartbeat для управления жизненным циклом ветки \texttt{/current}.
\end{itemize}

\begin{figure}[htbp]
\centering
\begin{adjustbox}{max width=\textwidth}
\begin{tikzpicture}[node distance=10mm and 16mm]
    \node[depbox] (sens) {DHT11 + MQ-2};
    \node[depbox, below=of sens] (fw) {Firmware loop()};
    \node[depbox, below left=of fw] (oled) {OLED SSD1306};
    \node[depbox, below=of fw] (net) {Wi-Fi + NTP};
    \node[depbox, below right=of fw] (fbw) {FirebaseClient\\(read/write)};
    \node[depgroup, fit=(sens)(fw)(oled)(net)(fbw), label=above:{\textbf{ESP8266 (course\_work)}}] {};

    \node[depbox, right=35mm of fw] (cur) {/current};
    \node[depbox, below=of cur] (his) {/history/\{epochHourStart\}};
    \node[depbox, below=of his] (sta) {/status/last\_viewed};
    \node[depgroup, fit=(cur)(his)(sta), label=above:{\textbf{Firebase RTDB}}] {};

    \node[depbox, right=35mm of his] (uc) {useCurrentData()};
    \node[depbox, above=of uc] (uh) {useHistoryData()};
    \node[depbox, below=of uc] (ui) {CurrentData + HistoryCharts};
    \node[depgroup, fit=(uc)(uh)(ui), label=above:{\textbf{React Frontend}}] {};

    \draw[deparrow] (sens) -- (fw);
    \draw[deparrow] (fw) -- (oled);
    \draw[deparrow] (fw) -- (net);
    \draw[deparrow] (fw) -- (fbw);
    \draw[deparrow] (fbw) -- (cur);
    \draw[deparrow] (fbw) -- (his);
    \draw[deparrow] (fbw) -- (sta);
    \draw[deparrow] (cur) -- (uc);
    \draw[deparrow] (his) -- (uh);
    \draw[deparrow] (uc) -- (ui);
    \draw[deparrow] (uh) -- (ui);
    \draw[deparrow] (uc.east) to[out=20,in=250] (sta.east);
\end{tikzpicture}

\end{adjustbox}
\caption{Диаграмма зависимостей подсистем}
\end{figure}

\subsection{Поток данных и управляющий поток}

В архитектуре одновременно работают два связанных потока:

\begin{enumerate}
    \item \textbf{Данные измерений:} датчики → прошивка → Firebase → frontend.
    \item \textbf{Сигнал активности пользователя:} frontend → Firebase (\texttt{/status/\allowbreak last\_viewed}) → прошивка.
\end{enumerate}

Такой подход обеспечивает баланс между актуальностью данных и экономией ресурсов: система публикует «живую» ветку только в тот момент, когда это действительно нужно пользователю.
